% GENERATED FILE. Edit canonical lessons, then run npm run generate:books.
% canonical-source-hash: fnv1a64:925ba937fc40106f
% canonical-lessons: KA-C113-asti, KA-C113-yajamana, KA-C113-sakti, KA-C113-ase, KA-C113-kanasu

\chapter{Property, Owner, Strength, Wish, Dream}
\label{ch:ka-property-owner-strength-wish-dream}
\hlchaptermodality{113}

\begin{chapteropening}
\textbf{By the end of this chapter:} \emph{I can say property, an owner, strength, a wish and a dream.}

Complete the last lesson of chapter 113: i can say property, an owner, strength, a wish and a dream.
\end{chapteropening}

\section[\texorpdfstring{\kn{ಆಸ್ತಿ} (āsti)}{ಆಸ್ತಿ (āsti)}]{\kn{ಆಸ್ತಿ} (āsti) — property}

\label{lesson:KA-C113-asti}



\begin{quote}
Before the new one: say the Kannada for a chance, then the Kannada for a tool.
\end{quote}

\subsection*{You'll want to know: \kn{ಆಸ್ತಿ}}
\textbf{\kn{ಆಸ್ತಿ}} — \emph{āsti} — \textquotedblleft{}property\textquotedblright{}.

\begin{grammarlens}[title={What you've built}]
One more word for this chapter.
\end{grammarlens}

\subsection*{Your turn}
\begin{itemize}
\raggedright
  \item \emph{Say it:} \emph{āsti}
  \item \emph{Say it:} \emph{āsti}, once more
  \item \emph{Say it:} say \emph{āsti}
  \item \emph{Recall:} say the Kannada for a chance, then the Kannada for a tool, then say \emph{āsti} again
\end{itemize}

\begin{tcolorbox}[breakable,colback=teal!4,colframe=teal!35!black,title={Before you move on}]
What does \kn{ಆಸ್ತಿ} mean? (\textquotedblleft{}Property\textquotedblright{}.) Say it once more.
\end{tcolorbox}
\section[\texorpdfstring{\kn{ಯಜಮಾನ} (yajamāna)}{ಯಜಮಾನ (yajamāna)}]{\kn{ಯಜಮಾನ} (yajamāna) — an owner, the master}

\label{lesson:KA-C113-yajamana}



\begin{quote}
Before the new one: say the Kannada for a tool, then the Kannada for property.
\end{quote}

\subsection*{You'll want to know: \kn{ಯಜಮಾನ}}
\textbf{\kn{ಯಜಮಾನ}} — \emph{yajamāna} — \textquotedblleft{}an owner, the master\textquotedblright{}.

\begin{grammarlens}[title={What you've built}]
One more word for this chapter.
\end{grammarlens}

\subsection*{Your turn}
\begin{itemize}
\raggedright
  \item \emph{Say it:} \emph{yajamāna}
  \item \emph{Say it:} \emph{yajamāna}, once more
  \item \emph{Say it:} say \emph{yajamāna}
  \item \emph{Recall:} say the Kannada for a tool, then the Kannada for property, then say \emph{yajamāna} again
\end{itemize}

\begin{tcolorbox}[breakable,colback=teal!4,colframe=teal!35!black,title={Before you move on}]
What does \kn{ಯಜಮಾನ} mean? (\textquotedblleft{}An owner, the master\textquotedblright{}.) Say it once more.
\end{tcolorbox}
\section[\texorpdfstring{\kn{ಶಕ್ತಿ} (śakti)}{ಶಕ್ತಿ (śakti)}]{\kn{ಶಕ್ತಿ} (śakti) — strength, power}

\label{lesson:KA-C113-sakti}



\begin{quote}
Before the new one: say the Kannada for property, then the Kannada for an owner.
\end{quote}

\subsection*{You'll want to know: \kn{ಶಕ್ತಿ}}
\textbf{\kn{ಶಕ್ತಿ}} — \emph{śakti} — \textquotedblleft{}strength, power\textquotedblright{}.

\begin{grammarlens}[title={What you've built}]
One more word for this chapter.
\end{grammarlens}

\subsection*{Your turn}
\begin{itemize}
\raggedright
  \item \emph{Say it:} \emph{śakti}
  \item \emph{Say it:} \emph{śakti}, once more
  \item \emph{Say it:} say \emph{śakti}
  \item \emph{Recall:} say the Kannada for property, then the Kannada for an owner, then say \emph{śakti} again
\end{itemize}

\begin{tcolorbox}[breakable,colback=teal!4,colframe=teal!35!black,title={Before you move on}]
What does \kn{ಶಕ್ತಿ} mean? (\textquotedblleft{}Strength, power\textquotedblright{}.) Say it once more.
\end{tcolorbox}
\section[\texorpdfstring{\kn{ಆಸೆ} (āse)}{ಆಸೆ (āse)}]{\kn{ಆಸೆ} (āse) — a wish, a desire}

\label{lesson:KA-C113-ase}



\begin{quote}
Before the new one: say the Kannada for an owner, then the Kannada for strength.
\end{quote}

\subsection*{You'll want to know: \kn{ಆಸೆ}}
\textbf{\kn{ಆಸೆ}} — \emph{āse} — \textquotedblleft{}a wish, a desire\textquotedblright{}.

\begin{grammarlens}[title={What you've built}]
One more word for this chapter.
\end{grammarlens}

\subsection*{Your turn}
\begin{itemize}
\raggedright
  \item \emph{Say it:} \emph{āse}
  \item \emph{Say it:} \emph{āse}, once more
  \item \emph{Say it:} say \emph{āse}
  \item \emph{Recall:} say the Kannada for an owner, then the Kannada for strength, then say \emph{āse} again
\end{itemize}

\begin{tcolorbox}[breakable,colback=teal!4,colframe=teal!35!black,title={Before you move on}]
What does \kn{ಆಸೆ} mean? (\textquotedblleft{}A wish, a desire\textquotedblright{}.) Say it once more.
\end{tcolorbox}
\section[\texorpdfstring{\kn{ಕನಸು} (kanasu)}{ಕನಸು (kanasu)}]{\kn{ಕನಸು} (kanasu) — a dream}

\label{lesson:KA-C113-kanasu}



\begin{quote}
Before the new one: say the Kannada for strength, then the Kannada for a wish.
\end{quote}

\subsection*{You'll want to know: \kn{ಕನಸು}}
\textbf{\kn{ಕನಸು}} — \emph{kanasu} — \textquotedblleft{}a dream\textquotedblright{}.

\begin{grammarlens}[title={What you've built}]
One more word for this chapter.
\end{grammarlens}

\subsection*{Your turn}
\begin{itemize}
\raggedright
  \item \emph{Say it:} \emph{kanasu}
  \item \emph{Say it:} \emph{kanasu}, once more
  \item \emph{Say it:} say \emph{kanasu}
  \item \emph{Recall:} say the Kannada for strength, then the Kannada for a wish, then say \emph{kanasu} again
\end{itemize}

\begin{tcolorbox}[breakable,colback=teal!4,colframe=teal!35!black,title={Before you move on}]
What does \kn{ಕನಸು} mean? (\textquotedblleft{}A dream\textquotedblright{}.) Say it once more.
\end{tcolorbox}
